
          %%%%% ~~~~~~~~~~~~~~~~~~~~ %%%%%

\begin{center}
\begin{minipage}{5.5in}
\begin{center}
\textbf{\large Summary}
\end{center}
This document outlines the methodology used by Statistics Canada’s Census of Environment 
(CoE) program to generate the area estimates by target geography and land cover type 
released on March 27, 2025. As had been advocated by a number of recent research 
articles \cite{Gallego2004,Olofsson201442,GreenCongalton2020,Stehman2023ErrorMatrix},
a design-based approach was adopted for the advantage that it yields nearly 
unbiased estimators as well as accompanying variance estimators. The sampling population 
is the collection of grid cells (100 metre $\times$ 100 metre tiles) resulting from partitioning the 
land mass of Canada according to one of the CoE reference geospatial grids. The sampling 
design was stratified simple random sampling, stratifying by the cross-tabulation of 
ecoprovinces and land cover types. The CoE Land Cover Register (LCR) Map Version 1.0 was 
used for stratification. Sample allocation decisions were guided by the objective of 
controlling the sampling uncertainty of key target variables, informed by the results of a 
suite of simulation studies. Reference data were generated by expert interpretation of
high-resolution satellite images of the selected grid cells/tiles acquired around the target period 
(calendar year 2020).
Horvitz-Thompson estimators and their associated variance estimators were used 
respectively to produce land cover area and associated variance estimates.
\end{minipage}
\end{center}

          %%%%% ~~~~~~~~~~~~~~~~~~~~ %%%%%
